% !TEX root = ../matan.tex

\section{Локальный экстремум функции нескольких переменных. Необходимое условие}

\subsection*{Определение локального экстремума}

\begin{Definition}[required]{Локальный экстремум}{}
	Пусть $f\colon\Omega\to\RR$, $\Omega\subset\RR^d$, и $x$ --- внутренняя точка
	области $\Omega$. Точка $x$ называется \textbf{точкой локального экстремума} функции
	$f$, если существует окрестность $U_x\subset\Omega$, в которой приращение
	$f(\xi)-f(x)$ сохраняет знак. А именно, для всех $\xi\in U_x$:
	\begin{itemize}
		\item $f(\xi)\le f(x)$ --- точка \emph{локального максимума}
		(строгого, если $f(\xi)<f(x)$ при $\xi\ne x$);
		\item $f(\xi)\ge f(x)$ --- точка \emph{локального минимума}
		(строгого, если $f(\xi)>f(x)$ при $\xi\ne x$).
	\end{itemize}
\end{Definition}

\subsection*{Необходимое условие экстремума}

\begin{Theorem}{Необходимое условие локального экстремума}{}
	Пусть $x$ --- точка локального экстремума функции $f\colon\Omega\to\RR$, и в этой
	точке существует градиент $\nabla f(x)$ (то есть существуют все частные
	производные). Тогда градиент обращается в нуль:
	\[
	\nabla f(x)=\bigl(\partial_1 f(x),\dots,\partial_d f(x)\bigr)=\vec 0.
	\]
\end{Theorem}

\begin{proof}
	Зафиксируем индекс $i\in\{1,\dots,d\}$ и все переменные, кроме $i$-й. Рассмотрим
	вспомогательную функцию одной переменной
	\[
	g_i(t)=f(x_1,\dots,x_{i-1},\,x_i+t,\,x_{i+1},\dots,x_d),
	\]
	определённую при достаточно малых $|t|$ (так что аргумент остаётся в $U_x$).

	Так как $x$ --- точка локального экстремума функции $f$, для всех малых $t$
	приращение $g_i(t)-g_i(0)=f(\dots,x_i+t,\dots)-f(x)$ сохраняет знак. Значит, $t=0$ ---
	точка локального экстремума функции $g_i$.

	Существование частной производной $\partial_i f(x)$ означает, что $g_i$
	дифференцируема в нуле, причём
	\[
	g_i'(0)=\lim_{t\to0}\frac{g_i(t)-g_i(0)}{t}=\partial_i f(x).
	\]
	По теореме Ферма (необходимое условие экстремума функции одной переменной)
	производная в точке внутреннего экстремума равна нулю, поэтому $g_i'(0)=0$, то есть
	\[
	\partial_i f(x)=0.
	\]
	Повторяя рассуждение для каждого $i=1,\dots,d$, заключаем, что все частные
	производные в точке $x$ равны нулю, а значит, $\nabla f(x)=\vec 0$.
\end{proof}

\begin{Definition}{Стационарная точка}{}
	Точка $x$, в которой $\nabla f(x)=\vec 0$, называется \textbf{стационарной точкой}
	функции $f$. Таким образом, необходимое условие можно сформулировать так:
	\emph{всякая точка локального экстремума, в которой существует градиент, является
	стационарной.}
\end{Definition}

\begin{Remark}{Условие не является достаточным}{}
	Обращение градиента в нуль не гарантирует экстремума. Классический пример --- 
	\emph{седловая точка}: для $f(x_1,x_2)=x_1^2-x_2^2$ в начале координат
	$\nabla f(0,0)=\vec 0$, однако вдоль оси $x_1$ функция возрастает
	($f(t,0)=t^2>0$), а вдоль оси $x_2$ убывает ($f(0,t)=-t^2<0$), так что в любой
	окрестности нуля есть точки и большего, и меньшего значения. Для различения
	экстремума и седла привлекают \emph{достаточные} условия (знакоопределённость
	второго дифференциала).
\end{Remark}

\begin{Remark}{Роль в поиске глобального экстремума}{}
	Необходимое условие даёт первый шаг стандартного алгоритма поиска глобального
	экстремума в области: сперва находят все стационарные точки внутри области (решая
	систему $\nabla f=\vec 0$), затем отдельно исследуют границу (задача на условный
	экстремум) и сравнивают значения $f$ во всех найденных точках.
\end{Remark}
